%% ma: L12-B08-0121
%% lop: 12
%% chuong: 2
%% bai: 8
%% dang: 12.8.5
%% loai: TN4
%% muc: TH
%% dap_an: "D"
%% nguon: "CĐ2.3.pdf giáo viên gửi 27/10/2026 (M07, Lê Minh Khôi – Nguyễn Anh Khôi, Chương II Chủ đề 3 Biểu thức tọa độ của các phép toán vectơ), Câu 33"
%% kien_thuc: "Bài 8 (hai vectơ cùng phương, ba điểm thẳng hàng)"
%% trang_thai: chinh-thuc
\begin{ex}
Trong không gian $Oxyz$, cho hai điểm $A(1;-2;3)$, $B(2;-3;4)$. Tọa độ của điểm $M\in(Oxy)$ sao cho ba điểm $A$, $B$, $M$ thẳng hàng là
\choice{$M(1;1;0)$}{$M(3;-5;7)$}{$M(-3;5;0)$}{\True $M(-2;1;0)$}
\loigiai{$M(x;y;0)$, $\overrightarrow{AB}=(1;-1;1)$, $\overrightarrow{AM}=(x-1;y+2;-3)=t\overrightarrow{AB}\Rightarrow t=-3$.

$x=-2$, $y=1$: $M(-2;1;0)$.}
\end{ex}
